% Part: propositional-logic
% Chapter: syntax-and-semantics
% Section: formulas

\documentclass[../../../include/open-logic-section]{subfiles}

\begin{document}

\olfileid{pl}{syn}{fml}

\olsection{Propositionele \usetoken{P}{formula}}

We bouwen de !!^{formula}s van de propositielogica op uit
\emph{!!{propositional
    variable}s}\iftag{prvFalse}{\iftag{prvTrue}{,}{ en} de
  propositionele constante~$\lfalse$}{}\iftag{prvTrue}{\iftag{prvFalse}{
    en}{} de propositionele constante~$\ltrue$}{}. Daarvoor gebruiken
we \emph{logische verbindingstekens}.

\begin{enumerate}
\item !!^a{denumerable} verzameling~$\PVar$ van !!{propositional variable}s
  $\Obj p_0$, $\Obj p_1$, \dots
\tagitem{prvFalse}{De propositionele constante voor !!{falsity}~$\lfalse$.}{}
\tagitem{prvTrue}{De propositionele constante voor !!{truth}~$\ltrue$.}{}
\item De logische verbindingstekens:
  \startycommalist
  \iftag{prvNot}{\ycomma $\lnot$ (negatie)}{}%
  \iftag{prvAnd}{\ycomma $\land$ (conjunctie)}{}%
  \iftag{prvOr}{\ycomma $\lor$ (disjunctie)}{}%
  \iftag{prvIf}{\ycomma $\lif$ (!!{conditional})}{}%
  \iftag{prvIff}{\ycomma $\liff$ (!!{biconditional})}{}%
\item Leestekens: (, ) en de komma.
\end{enumerate}

We noemen deze taal van de propositielogica $\Lang L_0$.

\iftag{defNot,defOr,defAnd,defIf,defIff,defTrue,defFalse,defEx,defAll}{%
Naast de primitieve verbindingstekens hierboven gebruiken we de
volgende \emph{gedefinieerde} symbolen:
\startycommalist
  \iftag{defNot}{\ycomma $\lnot$ (negatie)}{}%
  \iftag{defAnd}{\ycomma $\land$ (conjunctie)}{}%
  \iftag{defOr}{\ycomma $\lor$ (disjunctie)}{}%
  \iftag{defIf}{\ycomma $\lif$ (!!{conditional})}{}%
  \iftag{defIff}{\ycomma $\liff$ (!!{biconditional})}{}%
  \iftag{defFalse}{\ycomma $\lfalse$ (!!{falsity})}{}%
  \iftag{defTrue}{\ycomma $\ltrue$ (!!{truth})}}{}.

\begin{tagblock}{defNot,defOr,defAnd,defIf,defIff,defTrue,defFalse,defEx,defAll}
\begin{explain}
Een gedefinieerd symbool hoort officieel niet bij de taal. We voeren
het in als informele afkorting. Zonder zo'n afkorting kunnen formules
die alleen primitieve symbolen gebruiken behoorlijk lang worden. Een
afkorting maakt die formules dus korter. Nog belangrijker is dat ook
de bewijzen korter worden. Voor elk primitief symbool is namelijk een
afzonderlijk geval in een bewijs nodig. Meer primitieve symbolen
betekenen daarom langere bewijzen.
\end{explain}
\end{tagblock}

% Alternate symbols

\begin{intro}
Misschien ken je andere termen en symbolen dan de termen en symbolen
hierboven. Logicaboeken en docenten gebruiken voor ``negatie'' vaak
$\sim$, $\neg$ of~!\ en voor ``conjunctie'' vaak $\wedge$, $\cdot$ of
$\&$. Voor het ``conditioneel'' of de ``implicatie'' zie je vaak $\rightarrow$,
$\Rightarrow$ of $\supset$.
\iftag{prvIff,defIff}{Voor ``biconditioneel'', ``bi-implicatie'' of
  ``(materiële) equivalentie'' worden $\leftrightarrow$, $\Leftrightarrow$ en
  $\equiv$ gebruikt.}{}
\iftag{prvFalse,defFalse}{Het symbool $\lfalse$ heet ook wel
  ``onwaarheid'', ``falsum'', ``absurdum'' of ``bottom''.}{}
\iftag{prvTrue,defTrue}{Het symbool $\ltrue$ heet ook wel ``waarheid'',
  ``verum'' of ``top''.}{}
\end{intro}

\begin{defn}[Formule]
\ollabel{defn:formulas}
We definiëren de verzameling~$\Frm[L_0]$ van \emph{!!{formula}s} van
de propositielogica stap voor stap als volgt:
\begin{enumerate}
\tagitem{prvFalse}{$\lfalse$ is een atomaire !!{formula}.}{}

\tagitem{prvTrue}{$\ltrue$ is een atomaire !!{formula}.}{}

\item Elke !!{propositional variable}~$\Obj p_i$ is een atomaire
  !!{formula}.

\tagitem{prvNot}{Als $!A$ !!a{formula} is, dan is $\lnot !A$
  !!a{formula}.}{}

\tagitem{prvAnd}{Als $!A$ en $!B$ !!{formula}s zijn, dan is $(!A \land
  !B)$ !!a{formula}.}{}

\tagitem{prvOr}{Als $!A$ en $!B$ !!{formula}s zijn, dan is $(!A \lor !B)$
  !!a{formula}.}{}

\tagitem{prvIf}{Als $!A$ en $!B$ !!{formula}s zijn, dan is $(!A \lif !B)$
  !!a{formula}.}{}

\tagitem{prvIff}{Als $!A$ en $!B$ !!{formula}s zijn, dan is $(!A \liff !B)$
  !!a{formula}.}{}

\tagitem{limitClause}{Niets anders is !!a{formula}.}{}
\end{enumerate}
\end{defn}

\begin{explain}
Dit is een \emph{inductieve definitie} van !!{formula}s. Je kunt haar
zien als een bouwproces met oneindig veel stappen. In de eerste stap
verklaren we alle atomaire formules tot formules. Dat zijn de eerste
gevallen van de definitie: de gevallen voor
\iftag{prvTrue}{$\ltrue$, }{}%
\iftag{prvFalse}{$\lfalse$, }{}%
$\Obj p_i$. Een ``atomaire !!{formula}'' is hier dus een !!{formula}
met een van deze vormen.

De andere gevallen zijn regels waarmee we nieuwe !!{formula}s bouwen
uit !!{formula}s die al klaar zijn. In de tweede stap passen we die
regels toe op atomaire !!{formula}s. In de derde stap gebruiken we de
atomaire formules en de formules uit de tweede stap. Zo gaan we verder.
Iets is precies dan !!a{formula} als het uiteindelijk in een van deze
stappen wordt gebouwd.
\end{explain}

Stel dat we de formule $(!B \ast !C)$ uit $!B$ en $!C$ bouwen met een
verbindingsteken~$\ast$ voor twee plaatsen. Dan laten we het buitenste
paar haakjes vaak weg en schrijven we gewoon~$!B \ast !C$.

\begin{tagblock}{defNot,defOr,defAnd,defIf,defIff,defTrue,defFalse,defEx,defAll}
\begin{defn}
Voor formules met de gedefinieerde operatoren gelden de volgende
afkortingen:

\begin{tagenumerate}{defTrue,defFalse,defNot,defOr,defAnd,defIf,defIff,defEx,defAll}

\tagitem{defTrue}{$\ltrue$ is een afkorting voor
  \iftag{prvFalse}{$\lnot\lfalse$}{$(!A \lor \lnot !A)$ voor één vaste
    atomaire !!{formula}~$!A$}.}{}

\tagitem{defFalse}{$\lfalse$ is een afkorting voor
  \iftag{prvTrue}{$\lnot\ltrue$}{$(!A \land \lnot !A)$ voor één vaste
    atomaire !!{formula}~$!A$}.}{}

\tagitem{defNot}{$\lnot !A$ is een afkorting voor $!A \lif \lfalse$.}{}

\tagitem{defOr}{$!A \lor !B$ is een afkorting voor
  \iftag{prvAnd}{$\lnot(\lnot !A \land \lnot !B)$}{$\lnot !A \lif
    !B$}.}{}

\tagitem{defAnd}{$!A \land !B$ is een afkorting voor
  \iftag{prvOr}{$\lnot(\lnot !A \lor \lnot !B)$}{$\lnot (!A \lif
    \lnot !B)$}.}{}

\tagitem{defIf}{$!A \lif !B$ is een afkorting voor
  \iftag{prvOr}{$\lnot !A \lor !B$}{$\lnot (!A \land \lnot !B)$}.}{}

\tagitem{defIff}{$!A \liff !B$ is een afkorting voor $(!A \lif !B) \land (!B
  \lif !A)$.}{}
\end{tagenumerate}
\end{defn}
\end{tagblock}

\begin{defn}[Syntactische identiteit]
Het symbool $\ident$ zegt dat twee tekenreeksen syntactisch identiek
zijn. Dus $!A \ident !B$ geldt precies als $!A$ en $!B$ evenveel
symbolen hebben en op elke plaats hetzelfde symbool staat.
\end{defn}

Links en rechts van $\ident$ mogen ook tekenreeksen staan die aan
elkaar zijn gezet. Zo betekent $!A \ident (!B \lor !C)$ het volgende.
De tekenreeks~$!A$ is, in deze volgorde, opgebouwd uit een openend
haakje, de tekenreeks $!B$, het symbool $\lor$, de tekenreeks~$!C$ en
een sluitend haakje. Als deze identiteit geldt, weten we dus dat het
  eerste symbool van $!A$ een openend haakje is. We weten ook dat $!A$
  vanaf de tweede plaats de tekenreeks $!B$ bevat, dat daarna $\lor$
volgt, enzovoort.
\end{document}
